% Luc Destombe --- licence CC BY-NC-SA 4.0
% https://creativecommons.org/licenses/by-nc-sa/4.0/deed.fr
% Fichier autonome : compiler avec pdflatex, deux fois (signets, nombre de pages).
\def\lyceeversion{1.1}
\documentclass[11pt,a4paper]{article}
\makeatletter
\newif\ifeleve
\elevefalse

\newif\iflycee@cm
\newif\iflycee@palatino

\RequirePackage[utf8]{inputenc}
\RequirePackage[T1]{fontenc}
\RequirePackage[french]{babel}
\RequirePackage{etoolbox}

\newcommand{\ShorthandsOff}{\@ifundefined{shorthandoff}{}{\shorthandoff{;:!?}}}
\newcommand{\ShorthandsOn}{\@ifundefined{shorthandon}{}{\shorthandon{;:!?}}}
\ShorthandsOff
\AtBeginDocument{\ShorthandsOn}
\AtBeginEnvironment{tikzpicture}{\ShorthandsOff}

\RequirePackage{geometry}
\RequirePackage{fancyhdr}
\RequirePackage{lastpage}
\RequirePackage{multicol}
\RequirePackage{array}
\RequirePackage{enumitem}
\RequirePackage{xcolor}
\RequirePackage{needspace}

\RequirePackage{amsmath,amssymb}
\iflycee@cm\else\RequirePackage{fourier}\fi
\DeclareMathAlphabet{\mathcal}{OMS}{cmsy}{m}{n}
\renewcommand{\ttdefault}{lmtt}
\begingroup\catcode`\_=\active \catcode`\^=\active
\global\let_\sb \global\let^\sp
\endgroup
\AtBeginDocument{\catcode`\_=12 \mathcode`\_="8000
  \catcode`\^=12 \mathcode`\^="8000 }
\newlength{\retraitcalculs}
\setlength{\retraitcalculs}{2em}

\RequirePackage{tikz}
\usetikzlibrary{arrows.meta,calc,babel,shapes.misc}
\RequirePackage[most]{tcolorbox}
\tcbuselibrary{skins,breakable}

\definecolor{coulDef}{HTML}{1F4E79}
\definecolor{coulProp}{HTML}{7030A0}
\definecolor{coulRem}{HTML}{C55A11}
\definecolor{coulEx}{HTML}{2E7D32}
\definecolor{coulExo}{HTML}{B03A2E}
\definecolor{coulPart}{HTML}{1F4E79}
\definecolor{coulMeth}{HTML}{1B6E4F}
\definecolor{coulCourbe}{HTML}{0B5ED7}
\definecolor{coulCourbeB}{HTML}{C00000}
\definecolor{coulCourbeC}{HTML}{2E7D32}
\definecolor{coulGrille}{HTML}{8C8C8C}
\definecolor{coulRep}{HTML}{1B6E4F}
\definecolor{coulBar}{HTML}{2E7D32}

\newcommand{\R}{\mathbb{R}}
\newcommand{\Rs}{\mathbb{R}^{*}}
\newcommand{\N}{\mathbb{N}}
\newcommand{\Ns}{\mathbb{N}^{*}}
\newcommand{\Z}{\mathbb{Z}}
\newcommand{\Q}{\mathbb{Q}}
\renewcommand{\le}{\leqslant}
\renewcommand{\ge}{\geqslant}
\renewcommand{\approx}{\simeq}
\newcommand{\vect}[1]{\overrightarrow{#1}}
\newcommand{\norme}[1]{\left\lVert #1 \right\rVert}
\newcommand{\intervalle}[2]{\left[#1\,;#2\right]}
\RequirePackage{cancel}
\renewcommand{\CancelColor}{\color{coulExo}}
\newcommand{\fois}[1]{\times{\color{coulExo}#1}}
\newcommand{\barre}[1]{\cancel{#1}}
\newlength{\lycee@eqh}\newlength{\lycee@eqd}
\newsavebox{\lycee@eqa}\newsavebox{\lycee@eqb}\newsavebox{\lycee@eqc}
\newcommand{\lycee@eqcol}[1]{\hbox{$\renewcommand{\arraystretch}{1.35}%
  \begin{array}[t]{@{}>{\displaystyle}r@{}>{\displaystyle{}}l@{}}#1\end{array}$}}
\newcommand{\lycee@eqmesure}[1]{%
  \ifdim\ht#1>\lycee@eqh\setlength{\lycee@eqh}{\ht#1}\fi
  \ifdim\dp#1>\lycee@eqd\setlength{\lycee@eqd}{\dp#1}\fi}
\newcommand{\lycee@eqtrait}{\hspace{1.2em}\tikz[baseline=0pt]\draw[dash pattern=on 2pt off 2pt,
  line width=0.5pt] (0,\lycee@eqh) -- (0,-\lycee@eqd);\hspace{1.2em}}
\newcommand{\eqproduit}[2]{\par\smallskip\noindent C'est une équation produit :
  $A\times B=0 \iff A=0$ ou $B=0$.\par\nopagebreak\smallskip
  \begingroup\setlength{\lycee@eqh}{0pt}\setlength{\lycee@eqd}{0pt}%
  \setbox\lycee@eqa=\lycee@eqcol{#1}\setbox\lycee@eqb=\lycee@eqcol{#2}%
  \lycee@eqmesure\lycee@eqa\lycee@eqmesure\lycee@eqb%
  \noindent\hspace*{\retraitcalculs}\usebox\lycee@eqa\lycee@eqtrait\usebox\lycee@eqb\endgroup\par\smallskip}
\newcommand{\eqproduittrois}[3]{\par\smallskip\noindent C'est une équation produit :
  $A\times B\times C=0 \iff A=0$ ou $B=0$ ou $C=0$.\par\nopagebreak\smallskip
  \begingroup\setlength{\lycee@eqh}{0pt}\setlength{\lycee@eqd}{0pt}%
  \setbox\lycee@eqa=\lycee@eqcol{#1}\setbox\lycee@eqb=\lycee@eqcol{#2}\setbox\lycee@eqc=\lycee@eqcol{#3}%
  \lycee@eqmesure\lycee@eqa\lycee@eqmesure\lycee@eqb\lycee@eqmesure\lycee@eqc%
  \noindent\hspace*{\retraitcalculs}\usebox\lycee@eqa\lycee@eqtrait\usebox\lycee@eqb\lycee@eqtrait\usebox\lycee@eqc\endgroup
  \par\smallskip}

\tikzset{grille/.style={coulGrille,line width=0.4pt}}
\newcommand{\quadrillage}[4]{\draw[grille,step=1] (#1,#2) grid (#3,#4);}
\tikzset{
  croix/.style={cross out,draw,line width=0.6pt,inner sep=0pt,minimum size=4pt},
  croixdroite/.style={croix,rotate=45,line width=0.8pt},
}
\newcommand{\pt}[3]{\path (#1) node[croix]{} node[#2,font=\small,inner sep=2.5pt]{$#3$};}

\newcommand{\lycee@repere}[4]{%
  \draw[grille,step=1] (#1,#2) grid (#3,#4);
  \draw[-{Stealth[length=2mm]},thick] (#1,0)--({#3+0.4},0) node[below left=-1pt]{$x$};
  \draw[-{Stealth[length=2mm]},thick] (0,#2)--(0,{#4+0.4}) node[below left=-1pt]{$y$};
  \node[below left,font=\scriptsize,inner sep=1.5pt] at (0,0) {$0$};}
\newcommand{\lycee@gradx}[1]{\foreach \k/\l in {#1}{%
  \draw (\k,0.07)--(\k,-0.07) node[below,font=\scriptsize,inner sep=1.5pt]{$\l$};}}
\newcommand{\lycee@grady}[1]{\foreach \k/\l in {#1}{%
  \draw (0.07,\k)--(-0.07,\k) node[left,font=\scriptsize,inner sep=1.5pt]{$\l$};}}
\newcommand{\lycee@intff}[2]{\left[#1\,;#2\right]}
\newcommand{\lycee@intfo}[2]{\left[#1\,;#2\right[}
\newcommand{\lycee@intof}[2]{\left]#1\,;#2\right]}
\newcommand{\lycee@intoo}[2]{\left]#1\,;#2\right[}
\newcommand{\lycee@ens}[1]{\left\{#1\right\}}
\AtBeginDocument{%
  \providecommand{\repere}{\lycee@repere}%
  \providecommand{\gradx}{\lycee@gradx}%
  \providecommand{\grady}{\lycee@grady}%
  \providecommand{\Cf}{\mathcal{C}\sb{f}}%
  \providecommand{\Cg}{\mathcal{C}\sb{g}}%
  \providecommand{\intff}{\lycee@intff}%
  \providecommand{\intfo}{\lycee@intfo}%
  \providecommand{\intof}{\lycee@intof}%
  \providecommand{\intoo}{\lycee@intoo}%
  \providecommand{\ens}{\lycee@ens}}

\newlist{questions}{enumerate}{3}
\setlist[questions,1]{label=\arabic*\textdegree),leftmargin=*,itemsep=2pt,topsep=3pt}
\setlist[questions,2]{label=\alph*),leftmargin=*,itemsep=1pt,topsep=2pt}
\setlist[questions,3]{label=\roman*),leftmargin=*,itemsep=1pt,topsep=1pt}
\setlist[itemize]{label=\textbullet,leftmargin=*,itemsep=2pt,topsep=3pt}

\newcommand{\besoinplace}[1]{\par\penalty\@M\begingroup
  \dimen@=#1\relax\dimen@ii=\pagegoal\advance\dimen@ii-\pagetotal
  \ifdim\dimen@>\dimen@ii\ifdim\dimen@ii>\z@\vfil\fi\break\fi\endgroup}

\newcommand{\lycee@pied}{\thepage/\pageref{LastPage}}
\newcommand{\entete}[2]{%
  \pagestyle{fancy}\fancyhf{}%
  \fancyhead[L]{\footnotesize\color{coulPart}\textbf{#1}}%
  \fancyhead[R]{\footnotesize\color{coulPart}\textbf{#2}}%
  \fancyfoot[C]{\footnotesize\lycee@pied}%
  \renewcommand{\headrulewidth}{0.4pt}%
  \renewcommand{\headrule}{\hbox to\headwidth{%
    \color{coulPart!40}\leaders\hrule height \headrulewidth\hfill}}}

\RequirePackage{ccicons}
\newcommand{\lycee@licence}{{\scriptsize\color{gray}%
  \copyright\ Luc Destombe --- \ccbyncsa\ CC BY-NC-SA 4.0}}
\newif\iflycee@licence \lycee@licencetrue
\newcommand{\sanslicence}{\lycee@licencefalse}
\AtBeginDocument{\iflycee@licence\AddToHookNext{shipout/foreground}{%
  \put(0,0){\rlap{\kern\dimexpr1in+\hoffset+\oddsidemargin+\textwidth\relax
    \llap{\raisebox{-\dimexpr1in+\voffset+\topmargin+\headheight+\headsep
      +\textheight+\footskip\relax}{\lycee@licence}}}}}\fi}

\newcommand{\lycee@site}{faire-des-maths.fr}
\newcommand{\lycee@version}{{\scriptsize\color{gray}\lycee@site
  \ifdefined\lyceeversion\ --- Version~\lyceeversion\fi}}
\AtBeginDocument{\iflycee@licence\AddToHookNext{shipout/foreground}{%
  \put(0,0){\rlap{\kern\dimexpr1in+\hoffset+\oddsidemargin\relax
    \raisebox{-\dimexpr1in+\voffset+\topmargin+\headheight+\headsep
      +\textheight+\footskip\relax}{\lycee@version}}}}\fi}

\newcommand{\formulesserrees}{%
  \setlength{\abovedisplayskip}{3pt}\setlength{\belowdisplayskip}{3pt}%
  \setlength{\abovedisplayshortskip}{2pt}\setlength{\belowdisplayshortskip}{3pt}}

\setlength{\parindent}{0pt}

\RequirePackage[implicit=false,hidelinks,pdfencoding=auto,psdextra,bookmarksopen,
  bookmarksopenlevel=1,pdfcreator={},pdfproducer={}]{hyperref}
\RequirePackage{bookmark}
\hypersetup{pdfauthor={Luc Destombe},pdfkeywords={CC BY-NC-SA 4.0}}
\newcounter{lycee@signet}
\newcommand{\signet}[3][]{\stepcounter{lycee@signet}%
  \edef\lycee@dest{\if\relax\detokenize{#1}\relax
    signet.\arabic{lycee@signet}\else#1\fi}%
  \hypertarget{\lycee@dest}{}\bookmark[level=#2,dest=\lycee@dest]{#3}}
\pdfstringdefDisableCommands{%
  \def\R{R}\def\N{N}\def\Z{Z}\def\Q{Q}\def\vect#1{#1}\def\Cf{Cf}\def\Cg{Cg}\def\fois#1{×#1}%
  \def\barre#1{#1}\def\intervalle#1#2{[#1;#2]}\def\intff#1#2{[#1;#2]}%
  \def\textsuperscript#1{#1}\def\dfrac#1#2{#1/#2}\def\frac#1#2{#1/#2}%
  \def\vec#1{#1⃗}}%
\newcommand{\res}[1]{#1}
\geometry{top=2cm,bottom=1cm,left=1cm,right=1cm,headheight=14pt,footskip=0.6cm}

\RequirePackage{pgfplots}
\pgfplotsset{compat=1.18}
\RequirePackage{tkz-tab}

\newcounter{partiecours}
\renewcommand{\thepartiecours}{\Roman{partiecours}}
\newcommand{\partiecours}[1]{%
  \refstepcounter{partiecours}%
  \Needspace*{8\baselineskip}%
  \bigskip
  \noindent\begin{tcolorbox}[enhanced,colback=coulPart,colframe=coulPart,
      boxrule=0pt,arc=2pt,left=8pt,right=8pt,top=4pt,bottom=4pt,
      before skip=6pt,after skip=10pt]
    \color{white}\bfseries\large \thepartiecours{} --- #1%
    \signet[partie-\arabic{partiecours}]{1}{\thepartiecours{} --- #1}
  \end{tcolorbox}%
  \addcontentsline{toc}{section}{\thepartiecours{} --- #1}%
}

\newtcolorbox{boitecours}[3][]{%
  enhanced,breakable,
  colback=white,colframe=#2,coltitle=white,
  fonttitle=\bfseries,
  attach boxed title to top left={xshift=8pt,yshift=-\tcboxedtitleheight/2},
  boxed title style={colback=#2,arc=2pt,boxrule=0pt},
  boxrule=0.9pt,arc=2pt,
  left=8pt,right=8pt,top=8pt,bottom=6pt,
  title={#3},#1}

\newcounter{methode}
\newenvironment{definitionbox}[1][Définition]%
  {\begin{boitecours}[phantom={\signet{2}{#1}}]{coulDef}{#1}}{\end{boitecours}}
\newenvironment{proprietebox}[1][Propriété]%
  {\begin{boitecours}[phantom={\signet{2}{#1}}]{coulProp}{#1}}{\end{boitecours}}
\newenvironment{methodebox}[1][Méthode]{\stepcounter{methode}%
  \begin{boitecours}[phantom={\signet[methode-\arabic{methode}]{2}{#1}}]{coulMeth}{#1}}%
  {\end{boitecours}}

\newcommand{\etape}[1]{\par\smallskip\textbf{\color{coulMeth}#1}\par\nobreak\smallskip}
\newcommand{\enonce}[1]{\emph{#1}\par\smallskip}
\newcommand{\reponse}[1]{\par\smallskip
  {\footnotesize\itshape\color{black!55}Réponses : #1}}
\newcommand{\demo}[1]{\textbf{\color{coulMeth}Démonstration.} #1}
\newcommand{\afaire}[1]{%
  \par\smallskip
  \noindent{\small\color{coulExo}$\blacktriangleright$\ \textbf{À faire :}
    fiche d'exercices du chapitre, #1.}\par\smallskip}

\newenvironment{calculs}{\setlength{\jot}{5pt}\@fleqntrue\@mathmargin\retraitcalculs
  \setlength{\abovedisplayskip}{4pt}\setlength{\belowdisplayskip}{4pt}%
  \setlength{\abovedisplayshortskip}{4pt}\setlength{\belowdisplayshortskip}{4pt}%
  \csname alignat*\endcsname{2}}%
  {\csname endalignat*\endcsname}
\newcommand{\rem}[1]{\text{\small\itshape\color{black!60}#1}}
\newcommand{\explique}[2]{\par\smallskip\noindent
  \begin{minipage}[t]{0.57\linewidth}\raggedright #1\end{minipage}\hfill
  \begin{minipage}[t]{0.40\linewidth}\raggedright\leavevmode
    {\small\itshape\color{black!60}#2}\end{minipage}\par}
\newcommand{\cas}[1]{\par\smallskip\noindent\textbf{\color{coulExo}#1}\vadjust{\nobreak}}
\newcommand{\calc}[1]{\par\smallskip\textbf{\color{coulExo}#1}\par\nobreak}

\newtcolorbox{remarquebox}[1][Remarque]{%
  enhanced,breakable,colback=white,colframe=coulRem,
  boxrule=0pt,leftrule=2.5pt,arc=0pt,outer arc=0pt,
  left=8pt,right=6pt,top=5pt,bottom=5pt,
  before upper={\textbf{\color{coulRem}#1}\par\smallskip}}

\newtcolorbox{exemplebox}[1][Exemple]{%
  enhanced,breakable,colback=white,colframe=coulEx,
  boxrule=0pt,leftrule=2.5pt,arc=0pt,outer arc=0pt,
  left=8pt,right=6pt,top=5pt,bottom=5pt,
  before upper={\textbf{\color{coulEx}#1}\par\smallskip}}

\pgfplotsset{
  repere/.style={
    axis lines=middle,
    axis line style={-{Stealth[length=2.4mm]},thick},
    every axis x label/.style={at={(ticklabel* cs:1.0)},anchor=west},
    every axis y label/.style={at={(ticklabel* cs:1.0)},anchor=south},
    label style={font=\footnotesize},
    tick label style={font=\scriptsize},
    grid=both,
    major grid style={grille},
    minor tick num=0,
    tick align=inside,
    clip mode=individual,
  },
}
\tikzset{
  courbe/.style={coulCourbe,line width=1.1pt,smooth},
  courbeB/.style={coulCourbeB,line width=1.1pt,smooth},
  courbeC/.style={coulCourbeC,line width=1.1pt,smooth},
}

\newlist{etapes}{enumerate}{1}
\setlist[etapes]{label=\textbf{\arabic*.},leftmargin=*,itemsep=1pt,topsep=2pt}

\newlist{expressions}{enumerate}{1}
\setlist[expressions]{label={\Alph*\ $=$},leftmargin=*,itemsep=3pt,topsep=3pt}

\newcounter{exo}
\renewcommand{\theexo}{\ifnum\value{exo}<10 0\fi\arabic{exo}}
\newtcolorbox{exercice}[1][]{%
  enhanced,breakable,
  colback=white,colframe=coulExo,
  boxrule=0.9pt,arc=2pt,
  left=8pt,right=8pt,top=8pt,bottom=6pt,
  coltitle=white,fonttitle=\bfseries,
  attach boxed title to top left={xshift=8pt,yshift=-\tcboxedtitleheight/2},
  boxed title style={colback=coulExo,arc=2pt,boxrule=0pt},
  phantom={\signet{2}{Exercice \arabic{exo}}},
  title={Exercice \theexo\ifx\relax#1\relax\else{} --- #1\fi}}
\newenvironment{exo}[1][\relax]{\refstepcounter{exo}\begin{exercice}[#1]}{\end{exercice}}

  \newenvironment{correction}{\par}{\par}
  \newcommand{\autableau}{}
  \newcommand{\etapeexemples}{\etape{Exemples corrigés}}
  \newcommand{\etapetableau}[1]{\etape{#1}}
  \newenvironment{exempletableau}[1][Exemple]%
    {\begin{exemplebox}[#1]}{\end{exemplebox}}

\RequirePackage{comment}
\excludecomment{correctionavj}

\newcommand{\coursEntete}{}
\newcommand{\coursNiveau}{}
\newcommand{\coursTitre}{}
\newcommand{\coursSurTitre}{}
\newcommand{\coursSousTitre}{}

\newcommand{\enteteCours}[1]{\renewcommand{\coursEntete}{#1}}
\newcommand{\chapitrecours}[2]{\enteteCours{Chapitre #1 --- #2}}
\newcommand{\niveaucours}[1]{\renewcommand{\coursNiveau}{#1}}
\newcommand{\titrecours}[3][]{%
  \renewcommand{\coursTitre}{#2}%
  \renewcommand{\coursSurTitre}{#1}%
  \renewcommand{\coursSousTitre}{#3}}

\pagestyle{fancy}
\fancyhf{}
\fancyhead[L]{\footnotesize\color{coulPart}\textbf{\coursEntete}}
\fancyhead[R]{\footnotesize\color{coulPart}\textbf{\coursNiveau}}
\fancyfoot[C]{\footnotesize\thepage}
\renewcommand{\headrulewidth}{0.4pt}
\renewcommand{\headrule}{\hbox to\headwidth{\color{coulPart!40}\leaders\hrule height \headrulewidth\hfill}}

\setlength{\parskip}{2pt}

\newcommand{\lycee@mentioneleve}{\\[2pt]{\small\itshape Version élève}}
\AtBeginDocument{%
  \begin{center}
    \begin{tcolorbox}[enhanced,width=0.72\linewidth,colback=white,colframe=coulPart,
        boxrule=1.6pt,arc=3pt,halign=center,top=10pt,bottom=10pt]
      \ifx\coursSurTitre\empty
        {\Huge\bfseries\color{coulPart} \coursTitre}\\[4pt]
      \else
        {\Huge\bfseries\color{coulPart} \coursTitre}\\[3pt]
        {\Large\bfseries\color{coulPart} \coursSurTitre}\\[5pt]
      \fi
      {\normalsize \coursSousTitre}%
      
    \end{tcolorbox}
  \end{center}%
}

\RequirePackage{tikz}
\RequirePackage{xcolor}

\usetikzlibrary{decorations.text}
\definecolor{lycee@metal}{HTML}{1F4E79}
\definecolor{lycee@manche}{HTML}{8A5A2B}
\definecolor{lycee@etincelle}{HTML}{F28C28}
\newcommand{\lycee@pluriel}[2]{#1 #2\ifnum#1>1 s\fi}
\newcommand{\lycee@sceau}{%
  \edef\lycee@bas{\lycee@pluriel{\lyceerelectures}{relecture} et
    \lycee@pluriel{\lyceerodages}{rodage}}%
  \begin{tikzpicture}[line join=round]
    \fill[coulPart!10,draw=coulPart,line width=.8pt] (0:2.2)
      \foreach \i in {1,...,71}{ -- (\i*5:{isodd(\i)?2.1:2.2}) } -- cycle;
    \draw[coulPart,line width=1.2pt] (0,0) circle (1.93);
    \draw[coulPart,line width=.4pt] (0,0) circle (1.85);
    \draw[coulPart,line width=.4pt] (0,0) circle (1.2);
    \path[decorate,decoration={text along path,text align=center,
        text={|\fontsize{9.5}{11}\selectfont\bfseries\color{coulPart}|TEST{\'E} AVEC MES {\'E}L{\`E}VES}}]
      (180:1.45) arc (180:0:1.45);
    \path[decorate,decoration={text along path,text align=center,
        text/.expanded={|\noexpand\fontsize{8.5}{10}\noexpand\selectfont\noexpand\color{coulPart}|\lycee@bas}}]
      (195:1.68) arc (195:345:1.68);
    \begin{scope}[x=.25mm,y=-.25mm]
      \begin{scope}[shift={(-32,-40)}]%
        \fill[lycee@metal] (23,31) rectangle (41,49);
        \draw[lycee@etincelle,line width=.5pt,line cap=round]
          (20,30)--(15,27) (19,34)--(14,35) (44,30)--(49,27) (45,34)--(50,35);
        \fill[lycee@manche,rounded corners=.3mm] (38,24.6) rectangle (60,27.6);
        \fill[lycee@metal,rounded corners=.4mm] (25,20) rectangle (39,31);
        \fill[lycee@etincelle,rounded corners=.2mm] (25,20) rectangle (28,31);
      \end{scope}
    \end{scope}
    \ifdefined\lyceesceauplusieurs\def\lycee@annee{Ann{\'e}es~:}\else
      \def\lycee@annee{Ann{\'e}e}\fi
    \path[decorate,decoration={text along path,text align=center,
        text/.expanded={|\noexpand\fontsize{8}{10}\noexpand\selectfont\noexpand\color{coulPart}|%
          \unexpanded\expandafter{\lycee@annee} \unexpanded\expandafter{\lyceesceauannees}}}]
      (95:1.08) arc (95:445:1.08);
  \end{tikzpicture}}
\AtEndDocument{\ifdefined\lyceesceauannees
  \par\besoinplace{5cm}\vspace{1.5em}%
  \begin{center}\lycee@sceau\end{center}\fi}
\makeatother

\chapitrecours{9}{Primitives}
\niveaucours{1\textsuperscript{re} STI2D}
\titrecours{PRIMITIVES}{Cours --- Première STI2D}

\begin{document}

\newtcblisting{python}{%
  listing only,enhanced,breakable,
  colback=black!3,colframe=black!35,boxrule=0.5pt,arc=1pt,
  left=6pt,right=6pt,top=2pt,bottom=2pt,
  listing options={language=Python,basicstyle=\ttfamily\small,
    keywordstyle=\color{coulMeth}\bfseries,columns=fullflexible,
    keepspaces=true,showstringspaces=false,aboveskip=0pt,belowskip=0pt}}
\colorlet{coulEuler}{red!75!black}

\begin{remarquebox}[Comment utiliser ce cours]
Chaque partie commence par ce qu'il faut \textbf{savoir} (définitions, propriétés), puis
vient ce qu'il faut \textbf{savoir faire} : une méthode, avec des
\textbf{exemples corrigés} faits en classe, puis des exercices
\og\textbf{À vous de jouer}\fg{} à chercher seul.

\end{remarquebox}

\partiecours{Primitive d'une fonction}

\begin{minipage}[c]{0.58\linewidth}\raggedright
\textbf{But :} en physique, on connaît souvent la vitesse d'un mobile et on cherche sa
position ; or la vitesse est la dérivée de la position. Chercher une primitive, c'est
faire le chemin inverse de la dérivation : on connaît la dérivée, on cherche la
fonction.

\smallskip
Par exemple, en dérivant ${F(x)=x^{2}-5x+2}$, on trouve ${F'(x)=2x-5}$. On dit que $F$
est une \textbf{primitive} de ${f(x)=2x-5}$.
\end{minipage}\hfill
\begin{minipage}[c]{0.38\linewidth}\centering
\begin{tikzpicture}
  \node (F) at (0,0) {$F(x)=x^{2}-5x+2$};
  \node (f) at (0,-1.8) {$f(x)=2x-5$};
  \draw[-{Stealth[length=2mm]},thick] (F.south west) ++(0.6,0) to[bend right=35]
    node[left,font=\footnotesize]{on dérive} ($(f.north west)+(0.6,0)$);
  \draw[-{Stealth[length=2mm]},thick,coulMeth] (f.north east) ++(-0.6,0) to[bend right=35]
    node[right,font=\footnotesize,align=left]{on cherche\\une primitive} ($(F.south east)+(-0.6,0)$);
\end{tikzpicture}
\end{minipage}

\begin{definitionbox}[Définition 1 --- Primitive]
Une \textbf{primitive} d'une fonction $f$ est une fonction $F$ dont la dérivée est $f$ :
\[
  F'(x)=f(x).
\]
Dire que $F$ est une primitive de $f$, c'est donc dire que $f$ est la dérivée de $F$.
\end{definitionbox}

\begin{remarquebox}[Rappel --- Dérivées (chapitre 5)]
\begin{minipage}[c]{0.5\linewidth}\raggedright
Pour savoir si $F$ est une primitive de $f$, on dérive $F$ avec les formules du
chapitre 5, \textbf{terme par terme}, un nombre qui multiplie restant devant ; puis on
compare le résultat à $f(x)$.
\end{minipage}\hfill
\begin{minipage}[c]{0.46\linewidth}\centering
\renewcommand{\arraystretch}{2.4}\setlength{\tabcolsep}{12pt}
\begin{tabular}{|c|c|}
\hline
\textbf{Fonction} & \textbf{Dérivée}\\
\hline
$k$ (un nombre) & $0$\\
\hline
$x^{n}$ \quad ($n=1$, $2$, $3$\dots) & $n\,x^{n-1}$\\
\hline
$\dfrac{1}{x}$ \quad (pour $x\neq 0$) & $-\dfrac{1}{x^{2}}$\\
\hline
\end{tabular}
\end{minipage}
\end{remarquebox}

\begin{methodebox}[Méthode 1 --- Vérifier qu'une fonction est une primitive]
\etapeexemples
Dans chaque cas, dire si $F$ est une primitive de $f$.
\begin{questions}[label=\alph*)]
  \item $f(x)=6x-4$ \quad et \quad $F(x)=3x^{2}-4x+7$
  \item $f(x)=2x+\dfrac{3}{x^{2}}$ \quad et \quad $F(x)=x^{2}-\dfrac{3}{x}$ \quad (pour $x>0$)
  \item $f(x)=4x^{3}+1$ \quad et \quad $F(x)=x^{4}+x^{2}$
\end{questions}
\begin{correction}
\explique{\textbf{But :} on dérive $F$ et on compare $F'(x)$ à $f(x)$ ; s'ils sont égaux,
$F$ est une primitive de $f$.}{c'est la définition 1}
\cas{a)}
\begin{calculs}
F(x) &= 3x^{2}-4x+7 &\qquad& \rem{on part de $F$}\\
F'(x) &= 3\times 2x-4+0 && \rem{terme par terme : $x^{2}$ devient $2x$ ; $7$ devient $0$}\\
F'(x) &= 6x-4 &&
\end{calculs}
\explique{Conclusion : $F'(x)=f(x)$, donc \res{$F$ est une primitive de $f$}.}{la
constante $7$ a disparu en dérivant}
\cas{b)}
\begin{calculs}
F(x) &= x^{2}-\dfrac{3}{x} &\qquad& \rem{$\frac{3}{x}=3\times\frac{1}{x}$ : le $3$ reste
devant}\\
F'(x) &= 2x-3\times\left(-\dfrac{1}{x^{2}}\right) && \rem{la dérivée de $\frac{1}{x}$ est
$-\frac{1}{x^{2}}$}\\
F'(x) &= 2x+\dfrac{3}{x^{2}} && \rem{moins par moins : plus}
\end{calculs}
\explique{Conclusion : $F'(x)=f(x)$, donc \res{$F$ est une primitive de $f$}.}{on
retrouve exactement $f(x)$}
\cas{c)}
\begin{calculs}
F(x) &= x^{4}+x^{2} &\qquad& \rem{on part de $F$}\\
F'(x) &= 4x^{3}+2x && \rem{terme par terme}
\end{calculs}
\explique{Conclusion : $F'(x)\neq f(x)$, donc \res{$F$ n'est pas une primitive de
$f$}.}{on trouve $2x$ au lieu de $1$ : une différence suffit}
\end{correction}

\etape{À vous de jouer}
Dans chaque cas, dire si $F$ est une primitive de $f$.
\begin{questions}[label=\alph*)]
  \item $f(x)=8x+5$ \quad et \quad $F(x)=4x^{2}+5x-3$
  \item $f(x)=6x^{2}-\dfrac{2}{x^{2}}$ \quad et \quad $F(x)=2x^{3}+\dfrac{2}{x}$ \quad
        (pour $x>0$)
  \item $f(x)=3x^{2}-2$ \quad et \quad $F(x)=x^{3}-2x^{2}$
\end{questions}
\reponse{a) oui ;\; b)~oui ;\; c)~non, car $F'(x)=3x^{2}-4x$.}
\end{methodebox}

\afaire{exercices 1 à 4}

\partiecours{Toutes les primitives d'une fonction}

Les fonctions ${F(x)=3x^{2}-4x+7}$, ${G(x)=3x^{2}-4x}$ et ${H(x)=3x^{2}-4x-2}$ ont la
même dérivée, $6x-4$ : la constante ajoutée devient $0$ en dérivant. Une fonction a
donc une infinité de primitives.

\begin{proprietebox}[Propriété 1 --- Toutes les primitives]
\begin{minipage}[c]{0.6\linewidth}\raggedright
Si $F$ est une primitive de $f$, les primitives de $f$ sont les fonctions
\[
  F(x)+C, \quad\text{où $C$ est un nombre réel (une constante).}
\]
Deux primitives d'une même fonction diffèrent donc d'une constante (admis).

\smallskip
Sur le graphique, leurs courbes se déduisent l'une de l'autre en montant ou en
descendant ; au même $x$, elles ont la même pente. De même, deux trains qui roulent à
chaque instant à la même vitesse gardent toujours le même écart.

\smallskip
Une \textbf{condition initiale} ${F(x_{0})=y_{0}}$ (souvent la valeur au départ, en
$t=0$) fixe la constante $C$ : une seule primitive la vérifie.
\end{minipage}\hfill
\begin{minipage}[c]{0.36\linewidth}\centering
\begin{tikzpicture}[x=0.5cm,y=0.5cm]
  \repere{-3}{-3}{3}{5}
  \gradx{1} \grady{1}
  \draw[coulCourbeB,line width=1pt,domain=-2.4:2.4,samples=40] plot (\x,{0.5*\x*\x+2});
  \draw[coulCourbe,line width=1pt,domain=-2.4:2.4,samples=40] plot (\x,{0.5*\x*\x});
  \draw[coulCourbeC,line width=1pt,domain=-2.4:2.4,samples=40] plot (\x,{0.5*\x*\x-2});
  \draw[black!70,thick] (0.3,1.8)--(1.7,3.2);
  \draw[black!70,thick] (0.3,-0.2)--(1.7,1.2);
  \draw[black!70,thick] (0.3,-2.2)--(1.7,-0.8);
  \draw[-{Stealth[length=1.5mm]},thick,dashed] (-1.8,1.62)--(-1.8,3.58)
    node[midway,left,font=\footnotesize]{$+2$};
  \node[coulCourbeB,font=\footnotesize] at (-3.6,4.9) {$C=2$};
  \node[coulCourbe,font=\footnotesize] at (-3.6,0.6) {$C=0$};
  \node[coulCourbeC,font=\footnotesize] at (3.7,-1.2) {$C=-2$};
\end{tikzpicture}\par
{\footnotesize courbes de $F(x)+C$ ; en gris,\\ les tangentes en $x=1$}
\end{minipage}
\end{proprietebox}

\begin{methodebox}[Méthode 2 --- Trouver la primitive qui vérifie une condition initiale]
\etapeexemples
\begin{questions}[label=\alph*)]
  \item Soit la fonction $f$ définie sur $\R$ par $f(x)=6x-4$ ; la fonction ${F(x)=3x^{2}-4x+7}$ en est une primitive
        (méthode 1). Déterminer la primitive $G$ de $f$ telle que ${G(1)=5}$.
  \item Un drone d'inspection décolle du toit d'un bâtiment, à $12$~m du sol, et monte à
        la verticale. Sa vitesse (en m/s) au bout de $t$ secondes est ${v(t)=6-2t}$ ; une
        primitive de $v$ est ${F(t)=6t-t^{2}}$. Son altitude $h(t)$ (en m) est la
        primitive de $v$ telle que ${h(0)=12}$. Déterminer $h(t)$, puis l'altitude du
        drone au bout de $3$~s.
\end{questions}
\begin{correction}
\cas{a)}
\begin{calculs}
G(x) &= 3x^{2}-4x+C &\qquad& \rem{toutes les primitives : la constante $7$ devient $C$}\\
G(1) &= 3\times 1^{2}-4\times 1+C && \rem{on remplace $x$ par $1$}\\
G(1) &= -1+C &&
\end{calculs}
\explique{La condition ${G(1)=5}$ donne une équation d'inconnue $C$ :}{on la résout
comme au collège}
\begin{calculs}
-1+C &= 5 &\qquad& \\
C &= 6 && \rem{on ajoute $1$ des deux côtés}
\end{calculs}
\explique{Donc ${G(x)=\res{3x^{2}-4x+6}}$.}{contrôle : $G(1)=3-4+6=5$}
\cas{b)}
\begin{calculs}
h(t) &= 6t-t^{2}+C &\qquad& \rem{toutes les primitives de $v$}\\
h(0) &= 6\times 0-0^{2}+C && \rem{on remplace $t$ par $0$}\\
h(0) &= C &&
\end{calculs}
\explique{La condition ${h(0)=12}$ donne ${C=12}$, donc ${h(t)=\res{6t-t^{2}+12}}$.}{en
$t=0$, la constante est l'altitude de départ}
\begin{calculs}
h(3) &= 6\times 3-3^{2}+12 &\qquad& \rem{on remplace $t$ par $3$}\\
h(3) &= \res{21} &&
\end{calculs}
\explique{Au bout de $3$~s, le drone est à $21$~m du sol.}{on conclut par une phrase}
\end{correction}

\etape{À vous de jouer}
\begin{questions}[label=\alph*)]
  \item Soit la fonction $f$ définie sur $\R$ par $f(x)=8x+5$ ; la fonction ${F(x)=4x^{2}+5x-3}$ en est une primitive.
        Déterminer la primitive $G$ de $f$ telle que ${G(2)=10}$.
  \item Le drone d'un club d'aéromodélisme décolle d'une table, à $1$~m du sol, et monte
        à la verticale. Sa vitesse (en m/s) au bout de $t$ secondes est ${v(t)=8-2t}$ ;
        une primitive de $v$ est ${F(t)=8t-t^{2}}$. Son altitude $h(t)$ (en m) est la
        primitive de $v$ telle que ${h(0)=1}$. Déterminer $h(t)$, puis l'altitude du
        drone au bout de $4$~s.
\end{questions}
\reponse{a) $G(x)=4x^{2}+5x-16$ ;\; b)~$h(t)=8t-t^{2}+1$ et $h(4)=17$ : le drone est à
$17$~m du sol.}
\end{methodebox}

\afaire{exercices 5 à 8}

\partiecours{Primitive d'un polynôme}

En dérivant $x^{3}$, on obtient $3x^{2}$. Pour revenir à $x^{2}$, on fait le chemin
inverse : l'exposant augmente de $1$, puis on divise par ce nouvel exposant, car
${\left(\dfrac{x^{3}}{3}\right)'=\dfrac{3x^{2}}{3}}$, soit $x^{2}$.

\begin{proprietebox}[Propriété 2 --- Primitives des polynômes]
\begin{minipage}[c]{0.48\linewidth}\centering
\renewcommand{\arraystretch}{2.4}\setlength{\tabcolsep}{14pt}
\begin{tabular}{|c|c|}
\hline
\textbf{Fonction} $f(x)$ & \textbf{Une primitive} $F(x)$\\
\hline
$k$ (un nombre) & $kx$\\
\hline
$x$ & $\dfrac{x^{2}}{2}$\\
\hline
$x^{2}$ & $\dfrac{x^{3}}{3}$\\
\hline
$x^{3}$ & $\dfrac{x^{4}}{4}$\\
\hline
$x^{n}$ \quad ($n=4$, $5$, $6$\dots) & $\dfrac{x^{n+1}}{n+1}$\\
\hline
\end{tabular}
\end{minipage}\hfill
\begin{minipage}[c]{0.48\linewidth}\raggedright
On calcule \textbf{terme par terme}, et un nombre qui multiplie reste devant, comme
pour la dérivée : une primitive de $6x^{2}$ est ${6\times\dfrac{x^{3}}{3}=2x^{3}}$.

\smallskip
Pour contrôler, on dérive la primitive trouvée : on doit retrouver $f(x)$.
\end{minipage}
\end{proprietebox}

\begin{methodebox}[Méthode 3 --- Calculer une primitive d'un polynôme]
\etapeexemples
\begin{questions}[label=\alph*)]
  \item Déterminer une primitive $F$ sur $\R$ de ${f(x)=6x^{2}+4x-5}$.
  \item Déterminer une primitive $G$ sur $\R$ de ${g(x)=8x^{3}-3x^{2}+2x}$.
  \item Déterminer une primitive $H$ sur $\R$ de ${h(x)=x^{3}+5x-2}$.
  \item Déterminer la primitive $K$ sur $\R$ de ${k(x)=3x^{2}-2x+1}$ telle que ${K(1)=4}$.
\end{questions}
\begin{correction}
\cas{a)}
\begin{calculs}
F(x) &= 6\times\dfrac{x^{3}}{3}+4\times\dfrac{x^{2}}{2}-5x &\qquad& \rem{terme par
terme : $x^{2}$ donne $\frac{x^{3}}{3}$ ; $x$ donne $\frac{x^{2}}{2}$}\\
F(x) &= \res{2x^{3}+2x^{2}-5x} && \rem{on simplifie : $\frac{6}{3}=2$ et $\frac{4}{2}=2$}
\end{calculs}
\explique{Contrôle : $F'(x)=6x^{2}+4x-5$, c'est bien $f(x)$.}{on dérive pour vérifier}
\cas{b)}
\begin{calculs}
G(x) &= 8\times\dfrac{x^{4}}{4}-3\times\dfrac{x^{3}}{3}+2\times\dfrac{x^{2}}{2}
&\qquad& \rem{terme par terme : $x^{3}$ donne $\frac{x^{4}}{4}$}\\
G(x) &= \res{2x^{4}-x^{3}+x^{2}} && \rem{on simplifie les fractions}
\end{calculs}
\cas{c)}
\begin{calculs}
H(x) &= \dfrac{x^{4}}{4}+5\times\dfrac{x^{2}}{2}-2x &\qquad& \rem{$-2$ donne $-2x$}\\
H(x) &= \res{\dfrac{1}{4}x^{4}+\dfrac{5}{2}x^{2}-2x} && \rem{les fractions restent :
$\frac{5}{2}$ ne se simplifie pas}
\end{calculs}
\cas{d)}
\explique{Toutes les primitives de $k$ :}{partie II : on ajoute $C$}
\begin{calculs}
K(x) &= 3\times\dfrac{x^{3}}{3}-2\times\dfrac{x^{2}}{2}+x+C &\qquad& \rem{terme par
terme}\\
K(x) &= x^{3}-x^{2}+x+C &&
\end{calculs}
\explique{La condition ${K(1)=4}$ :}{on remplace $x$ par $1$}
\begin{calculs}
K(1) &= 1^{3}-1^{2}+1+C &\qquad& \\
K(1) &= 1+C &&
\end{calculs}
\begin{calculs}
1+C &= 4 &\qquad& \rem{$K(1)$ doit valoir $4$}\\
C &= 3 && \rem{on retire $1$ des deux côtés}
\end{calculs}
\explique{Donc ${K(x)=\res{x^{3}-x^{2}+x+3}}$.}{contrôle : $K(1)=1-1+1+3=4$}
\end{correction}

\etape{À vous de jouer}
\begin{questions}[label=\alph*)]
  \item Déterminer une primitive $F$ sur $\R$ de ${f(x)=9x^{2}-6x+2}$.
  \item Déterminer une primitive $G$ sur $\R$ de ${g(x)=4x^{3}+6x^{2}-8x}$.
  \item Déterminer une primitive $H$ sur $\R$ de ${h(x)=x^{3}-7x+1}$.
  \item Déterminer la primitive $K$ sur $\R$ de ${k(x)=6x^{2}+2x-3}$ telle que ${K(2)=10}$.
\end{questions}
\reponse{a) $F(x)=3x^{3}-3x^{2}+2x$ ;\; b)~$G(x)=x^{4}+2x^{3}-4x^{2}$ ;\;
c)~$H(x)=\frac{1}{4}x^{4}-\frac{7}{2}x^{2}+x$ ;\; d)~$K(x)=2x^{3}+x^{2}-3x-4$.}
\end{methodebox}

\afaire{exercices 9 à 13}

\partiecours{Primitives de cosinus et sinus}

Au chapitre 5, la dérivée de $\sin x$ est $\cos x$, et celle de $\cos x$ est $-\sin x$
($x$ est un angle en radians). Donc $\sin x$ est une primitive de $\cos x$. De même,
${(-\cos x)'=-(-\sin x)}$, soit $\sin x$ : $-\cos x$ est une primitive de $\sin x$.

\begin{proprietebox}[Propriété 3 --- Primitives de cosinus et sinus]
\begin{minipage}[c]{0.6\linewidth}\centering
\renewcommand{\arraystretch}{1.8}\setlength{\tabcolsep}{14pt}
\begin{tabular}{|c|c|}
\hline
\textbf{Fonction} $f(x)$ & \textbf{Une primitive} $F(x)$\\
\hline
$\cos x$ & $\sin x$\\
\hline
$\sin x$ & $-\cos x$\\
\hline
\end{tabular}

\smallskip\raggedright
Attention au signe $-$ : seule la primitive de $\sin x$ en a un. On calcule terme par
terme, et un nombre qui multiplie reste devant, comme pour les polynômes.
\end{minipage}\hfill
\begin{minipage}[c]{0.36\linewidth}\centering
\begin{tikzpicture}[font=\small]
  \node (s) at (0,1.5) {$\sin x$};
  \node (c) at (2.6,1.5) {$\cos x$};
  \node (ms) at (2.6,0) {$-\sin x$};
  \node (mc) at (0,0) {$-\cos x$};
  \draw[-{Stealth[length=2mm]},thick] (s)--(c);
  \draw[-{Stealth[length=2mm]},thick] (c)--(ms);
  \draw[-{Stealth[length=2mm]},thick] (ms)--(mc);
  \draw[-{Stealth[length=2mm]},thick] (mc)--(s);
\end{tikzpicture}\par
{\footnotesize on dérive : sens des flèches ;\\ une primitive : sens contraire}
\end{minipage}
\end{proprietebox}

\begin{methodebox}[Méthode 4 --- Calculer une primitive avec cosinus et sinus]
\etapeexemples
\begin{questions}[label=\alph*)]
  \item Déterminer une primitive $F$ sur $\R$ de ${f(x)=3\cos x}$.
  \item Déterminer une primitive $G$ sur $\R$ de ${g(x)=2x-5\sin x}$.
  \item Déterminer la primitive $H$ sur $\R$ de ${h(x)=4\cos x+\sin x}$ telle que ${H(0)=2}$.
\end{questions}
\begin{correction}
\cas{a)}
\begin{calculs}
F(x) &= \res{3\sin x} &\qquad& \rem{$\cos x$ donne $\sin x$ ; le $3$ reste devant}
\end{calculs}
\cas{b)}
\begin{calculs}
G(x) &= 2\times\dfrac{x^{2}}{2}-5\times(-\cos x) &\qquad& \rem{$x$ donne
$\frac{x^{2}}{2}$ ; $\sin x$ donne $-\cos x$}\\
G(x) &= \res{x^{2}+5\cos x} && \rem{moins par moins : plus}
\end{calculs}
\cas{c)}
\explique{Toutes les primitives de $h$ :}{on ajoute $C$}
\begin{calculs}
H(x) &= 4\sin x+(-\cos x)+C &\qquad& \rem{terme par terme}\\
H(x) &= 4\sin x-\cos x+C &&
\end{calculs}
\explique{La condition ${H(0)=2}$ :}{on remplace $x$ par $0$}
\begin{calculs}
H(0) &= 4\sin 0-\cos 0+C &\qquad& \\
H(0) &= 4\times 0-1+C && \rem{$\sin 0=0$ et $\cos 0=1$}\\
H(0) &= -1+C &&
\end{calculs}
\begin{calculs}
-1+C &= 2 &\qquad& \rem{$H(0)$ doit valoir $2$}\\
C &= 3 && \rem{on ajoute $1$ des deux côtés}
\end{calculs}
\explique{Donc ${H(x)=\res{4\sin x-\cos x+3}}$.}{contrôle : $H(0)=0-1+3=2$}
\end{correction}

\etape{À vous de jouer}
\begin{questions}[label=\alph*)]
  \item Déterminer une primitive $F$ sur $\R$ de ${f(x)=-2\cos x}$.
  \item Déterminer une primitive $G$ sur $\R$ de ${g(x)=4x+3\sin x}$.
  \item Déterminer la primitive $H$ sur $\R$ de ${h(x)=2\cos x-3\sin x}$ telle que ${H(0)=5}$.
\end{questions}
\reponse{a) $F(x)=-2\sin x$ ;\; b)~$G(x)=2x^{2}-3\cos x$ ;\;
c)~$H(x)=2\sin x+3\cos x+2$.}
\end{methodebox}

\afaire{exercices 14 à 17}

\partiecours{Primitives de $A\cos(\omega t+\varphi)$ et $A\sin(\omega t+\varphi)$}

En électricité et en mécanique, une grandeur qui oscille s'écrit $A\cos(\omega t+\varphi)$
ou $A\sin(\omega t+\varphi)$ (comme au chapitre 7, $A$ est l'amplitude, $\omega$ la
pulsation, $\varphi$ la phase à l'origine). Au chapitre 5, la dérivée de $\sin(3t)$ est
$3\cos(3t)$ : le $3$ passe devant. Pour une primitive de $\cos(3t)$, on divise donc par
$3$, et on vérifie en dérivant :
${\left(\dfrac{1}{3}\sin(3t)\right)'=\dfrac{1}{3}\times 3\cos(3t)=\cos(3t)}$.

\begin{proprietebox}[Propriété 4 --- Primitives de $A\cos(\omega t+\varphi)$ et $A\sin(\omega t+\varphi)$]
\begin{center}
\renewcommand{\arraystretch}{2.4}\setlength{\tabcolsep}{14pt}
\begin{tabular}{|c|c|}
\hline
\textbf{Fonction} $f(t)$ & \textbf{Une primitive} $F(t)$\\
\hline
$A\cos(\omega t+\varphi)$ & $\dfrac{A}{\omega}\sin(\omega t+\varphi)$\\
\hline
$A\sin(\omega t+\varphi)$ & $-\dfrac{A}{\omega}\cos(\omega t+\varphi)$\\
\hline
\end{tabular}
\end{center}
On divise par la pulsation $\omega$, le nombre devant $t$ ; la parenthèse ne change pas.
Comme en partie IV, seule la primitive du sinus a un signe $-$.
\end{proprietebox}

\begin{methodebox}[Méthode 5 --- Calculer une primitive de $A\cos(\omega t+\varphi)$ ou $A\sin(\omega t+\varphi)$]
\etapeexemples
\begin{questions}[label=\alph*)]
  \item Déterminer une primitive $F$ sur $\R$ de ${f(t)=6\cos(2t)}$.
  \item Déterminer une primitive $G$ sur $\R$ de ${g(t)=10\sin\left(5t+\dfrac{\pi}{4}\right)}$.
  \item Une masse accrochée à un ressort oscille de haut en bas. Sa vitesse (en cm/s) au
        bout de $t$ secondes est ${v(t)=6\sin(2t)}$, et sa hauteur au-dessus de la table
        (en cm) est la primitive $h$ de $v$ telle que ${h(0)=10}$. Déterminer $h(t)$,
        puis la hauteur de la masse à l'instant ${t=\dfrac{\pi}{2}}$.
\end{questions}
\begin{correction}
\cas{a)}
\begin{calculs}
F(t) &= \dfrac{6}{2}\sin(2t) &\qquad& \rem{$A=6$ et $\omega=2$ : on divise par $2$}\\
F(t) &= \res{3\sin(2t)} &&
\end{calculs}
\explique{Contrôle : $F'(t)=3\times 2\cos(2t)$, soit $6\cos(2t)$.}{on dérive pour
vérifier}
\cas{b)}
\begin{calculs}
G(t) &= -\dfrac{10}{5}\cos\left(5t+\dfrac{\pi}{4}\right) &\qquad& \rem{$A=10$ et
$\omega=5$ ; sinus : signe $-$}\\
G(t) &= \res{-2\cos\left(5t+\dfrac{\pi}{4}\right)} && \rem{la parenthèse ne change pas}
\end{calculs}
\cas{c)}
\explique{Toutes les primitives de $v$ :}{on ajoute $C$}
\begin{calculs}
h(t) &= -\dfrac{6}{2}\cos(2t)+C &\qquad& \rem{$A=6$ et $\omega=2$ ; sinus : signe $-$}\\
h(t) &= -3\cos(2t)+C &&
\end{calculs}
\explique{La condition ${h(0)=10}$ :}{on remplace $t$ par $0$}
\begin{calculs}
h(0) &= -3\cos(2\times 0)+C &\qquad& \\
h(0) &= -3+C && \rem{$\cos 0=1$}
\end{calculs}
\begin{calculs}
-3+C &= 10 &\qquad& \rem{$h(0)$ doit valoir $10$}\\
C &= 13 && \rem{on ajoute $3$ des deux côtés}
\end{calculs}
\explique{Donc ${h(t)=\res{13-3\cos(2t)}}$.}{contrôle : $h(0)=13-3=10$}
\begin{calculs}
h\left(\dfrac{\pi}{2}\right) &= 13-3\cos\left(2\times\dfrac{\pi}{2}\right) &\qquad&
\rem{on remplace $t$ par $\frac{\pi}{2}$}\\
h\left(\dfrac{\pi}{2}\right) &= 13-3\times(-1) && \rem{$\cos\pi=-1$}\\
h\left(\dfrac{\pi}{2}\right) &= \res{16} &&
\end{calculs}
\explique{À l'instant $t=\dfrac{\pi}{2}$ (environ $1{,}6$~s), la masse est à $16$~cm de la
table.}{c'est sa position la plus haute}
\end{correction}

\etape{À vous de jouer}
\begin{questions}[label=\alph*)]
  \item Déterminer une primitive $F$ sur $\R$ de ${f(t)=12\cos(4t)}$.
  \item Déterminer une primitive $G$ sur $\R$ de ${g(t)=6\sin\left(3t-\dfrac{\pi}{6}\right)}$.
  \item Une autre masse, accrochée à un ressort plus raide, oscille de haut en bas. Sa
        vitesse (en cm/s) au bout de $t$ secondes est ${v(t)=10\sin(5t)}$, et sa hauteur
        au-dessus de la table (en cm) est la primitive $h$ de $v$ telle que ${h(0)=20}$.
        Déterminer $h(t)$, puis la hauteur de la masse à l'instant ${t=\dfrac{\pi}{5}}$.
\end{questions}
\reponse{a) $F(t)=3\sin(4t)$ ;\; b)~$G(t)=-2\cos\left(3t-\frac{\pi}{6}\right)$ ;\;
c)~$h(t)=22-2\cos(5t)$ et $h\left(\frac{\pi}{5}\right)=24$ : la masse est à $24$~cm de la
table.}
\end{methodebox}

\afaire{exercices 18 à 22}

\partiecours{Méthode d'Euler}

\textbf{But :} aucune formule du cours ne donne une primitive de ${f(t)=\dfrac{1}{t}}$.
On veut quand même connaître la primitive $F$ de $f$ qui vérifie une condition
initiale, par exemple ${F(1)=0}$ : la méthode d'Euler en calcule des valeurs
approchées, de proche en proche, et trace sa courbe point par point. En physique, ce
problème s'écrit ${y'=f(t)}$ et ${y(t_{0})=y_{0}}$ : on cherche $y$, connaissant sa
dérivée et sa valeur au départ.

\smallskip
Le lien avec le chapitre 5 : près de $t$, la courbe de $F$ est presque confondue avec
sa tangente (approximation affine). Pour un petit pas $h$, on a donc
${F(t+h)\approx F(t)+h\times F'(t)}$, et ici ${F'(t)=f(t)}$.

\begin{proprietebox}[Propriété 5 --- Méthode d'Euler]
\begin{minipage}[c]{0.58\linewidth}\raggedright
On part de ${F(t_{0})=y_{0}}$ et on choisit un pas $h$ petit. On avance de $h$ en suivant
la tangente :
\[
  F(t+h)\approx F(t)+h\times f(t).
\]
En mots : \textbf{nouvelle valeur ${}\approx{}$ ancienne valeur ${}+h\times f(\text{ancien }t)$}. On recommence à partir du point obtenu. Les points sont proches de la courbe
de $F$, et d'autant plus que le pas $h$ est petit.
\end{minipage}\hfill
\begin{minipage}[c]{0.38\linewidth}\centering
\begin{tikzpicture}[x=1cm,y=1cm]
  \repere{0}{0}{4}{2}
  \gradx{1,2,3} \grady{1}
  \draw[coulCourbe,line width=1.1pt,domain=1:4.2,samples=40] plot (\x,{ln(\x)});
  \draw[coulEuler,thick] (1,0)--(2,1)--(3,1.5)--(4,1.833);
  \path[coulEuler] (1,0) node[croix]{} (2,1) node[croix]{} (3,1.5) node[croix]{}
    (4,1.833) node[croix]{};
  \node[coulCourbe,font=\footnotesize] at (3.7,0.95) {$\mathcal{C}_{F}$};
\end{tikzpicture}\par
{\footnotesize en rouge : la méthode d'Euler\\ avec un pas $h=1$, pour $f(t)=\frac{1}{t}$}
\end{minipage}
\end{proprietebox}

\begin{methodebox}[Méthode 6 --- Construire point par point une approximation d'une primitive]
\etapeexemples
Soit ${f(t)=\dfrac{1}{t}}$ pour $t>0$, et $F$ la primitive de $f$ telle que ${F(1)=0}$.
\begin{questions}[label=\alph*)]
  \item Avec la méthode d'Euler et un pas ${h=0{,}5}$, calculer des valeurs approchées
        de $F(1{,}5)$ ; $F(2)$ ; $F(2{,}5)$ et $F(3)$, arrondies à $0{,}01$.
  \item Placer les points obtenus dans un repère et les relier : c'est une approximation
        de la courbe de $F$ sur $\intff{1}{3}$.
\end{questions}
\begin{correction}
\cas{a)}
\explique{On part de ${t=1}$, où ${F(1)=0}$, et on avance de $0{,}5$ à chaque étape,
jusqu'à ${t=3}$ : quatre étapes.}{la condition initiale donne le premier point}
\begin{calculs}
F(1{,}5) &\approx F(1)+0{,}5\times f(1) &\qquad& \rem{nouvelle valeur ${}\approx{}$
ancienne ${}+h\times f(\text{ancien } t)$}\\
F(1{,}5) &\approx 0+0{,}5\times 1 && \rem{$F(1)=0$ et $f(1)=\frac{1}{1}=1$}\\
F(1{,}5) &\approx \res{0{,}5} &&
\end{calculs}
\explique{On recommence à partir de ${t=1{,}5}$, une étape par ligne :}{la valeur
trouvée devient l'ancienne valeur}
\begin{calculs}
F(2) &\approx 0{,}5+0{,}5\times\dfrac{1}{1{,}5}\approx\res{0{,}83} &\qquad&
\rem{$f(1{,}5)=\frac{1}{1{,}5}$}\\
F(2{,}5) &\approx 0{,}83+0{,}5\times\dfrac{1}{2}=\res{1{,}08} && \rem{$f(2)=\frac{1}{2}$}\\
F(3) &\approx 1{,}08+0{,}5\times\dfrac{1}{2{,}5}=\res{1{,}28} &&
\rem{$f(2{,}5)=\frac{1}{2{,}5}$}
\end{calculs}
\cas{b)}
\par\nopagebreak
\begin{minipage}[c]{0.44\linewidth}\raggedright
\renewcommand{\arraystretch}{1.5}\setlength{\tabcolsep}{4pt}
\begin{tabular}{|c|c|c|c|c|c|}
\hline
$t$ & $1$ & $1{,}5$ & $2$ & $2{,}5$ & $3$\\
\hline
$F(t)\approx$ & $0$ & $0{,}5$ & $0{,}83$ & $1{,}08$ & $1{,}28$\\
\hline
\end{tabular}

\smallskip
On place les points du tableau, puis on les relie par des segments : ce sont des
morceaux de tangentes.
\end{minipage}\hfill
\begin{minipage}[c]{0.52\linewidth}\centering
\begin{tikzpicture}[x=0.75cm,y=0.75cm]
  \quadrillage{0}{0}{7}{3}
  \draw[-{Stealth[length=2mm]},thick] (0,0)--(7.4,0) node[below left=-1pt]{$t$};
  \draw[-{Stealth[length=2mm]},thick] (0,0)--(0,3.4) node[below left=-1pt]{$y$};
  \node[below left,font=\scriptsize,inner sep=1.5pt] at (0,0) {$0$};
  \foreach \k/\l in {2/1,4/2,6/3}{\draw (\k,0.07)--(\k,-0.07)
    node[below,font=\scriptsize,inner sep=1.5pt]{$\l$};}
  \foreach \k/\l in {1/{0{,}5},2/1}{\draw (0.07,\k)--(-0.07,\k)
    node[left,font=\scriptsize,inner sep=1.5pt]{$\l$};}
  \draw[coulEuler,thick] (2,0)--(3,1)--(4,1.66)--(5,2.16)--(6,2.56);
  \path[coulEuler] (2,0) node[croix]{} (3,1) node[croix]{} (4,1.66) node[croix]{}
    (5,2.16) node[croix]{} (6,2.56) node[croix]{};
\end{tikzpicture}
\end{minipage}
\end{correction}

\etape{À vous de jouer}
Soit ${g(t)=\dfrac{1}{1+t^{2}}}$, et $G$ la primitive de $g$ telle que ${G(0)=0}$.
\begin{questions}[label=\alph*)]
  \item Avec la méthode d'Euler et un pas ${h=0{,}5}$, calculer des valeurs approchées
        de $G(0{,}5)$ ; $G(1)$ ; $G(1{,}5)$ et $G(2)$, arrondies à $0{,}01$.
  \item Placer les points obtenus dans un repère et les relier : c'est une approximation
        de la courbe de $G$ sur $\intff{0}{2}$.
\end{questions}
\reponse{a) $G(0{,}5)\approx 0{,}5$ ;\; $G(1)\approx 0{,}9$ ;\; $G(1{,}5)\approx 1{,}15$
;\; $G(2)\approx 1{,}30$.}
\end{methodebox}

\begin{remarquebox}[Avec un programme]
Plus le pas est petit, plus les points sont proches de la courbe de $F$, mais plus il y
a d'étapes : on les confie à un programme. Les deux lignes de la boucle sont celles de
la méthode : \texttt{y} augmente de $h\times f(t)$, puis \texttt{t} avance de $h$.

\smallskip
\noindent
\begin{minipage}[t]{0.42\linewidth}\vspace{0pt}
\begin{python}
def f(t):
    return 1/t

def euler(t, y, h, n):
    for i in range(n):
        y = y + h*f(t)
        t = t + h
    return y
\end{python}
\end{minipage}\hfill
\begin{minipage}[t]{0.54\linewidth}\vspace{0pt}\raggedright
\texttt{euler(1, 0, 0.5, 4)} fait les quatre étapes de la méthode 6 et renvoie
\texttt{1.2833...}, la valeur approchée de $F(3)$.

\smallskip
Avec un pas de $0{,}01$, \texttt{euler(1, 0, 0.01, 200)} renvoie \texttt{1.1019...} :
la valeur de $F(3)$ est plus précise.

\smallskip
Pour une autre fonction, on change seulement la ligne \texttt{return}.
\end{minipage}
\end{remarquebox}

\afaire{exercices 23 à 26}

\end{document}
